\section{Egenværdier}
\subsection{Definition}
For vektorrum $V$ på legeme $\mathbb{F}$ og lineær afbildning $L: V \to V$ 
er en skalar $\lambda \in \mathbb{F}$ en \textbf{egenværdi} for $L$, hvis der findes en vektor $\mathbf{v} \in V$, $\mathbf{v} \neq 0$, sådan at
\[L(v) = \lambda \cdot \mathbf{v}.\]
Vektoren $\mathbf{v}$ er en \textbf{egenvektor} tilhørende egenværdien $\lambda$.

\bemærk{$\lambda$ må godt være $0$}
\kilde{Definition 12.1.1}

\subsubsection{Matricers egenværdier}
For en matrix $A \in \mathbb{F}^{n \times n}$ og en vektor $\mathbf{v} \in \mathbb{F}^n$, $\mathbf{v} \neq 0$, er en skalar $\lambda \in \mathbb{F}$ en \textbf{egenværdi} for $A$, hvis
\[A \mathbf{v} = \lambda \cdot \mathbf{v}.\]
Vektoren $\mathbf{v}$ er en \textbf{egenvektor} tilhørende egenværdien $\lambda$.

\bemærk{Gælder kun for kvadratiske matricer}
\kilde{Definition 12.1.2}

\subsection{Løs Egenværdiproblemet}
For matrix $A \in \mathbb{F}^{n \times n}$, vil en værdi $\lambda \in \mathbb{F}$ være en egenværdi for $A$ hvis og kun hvis
\[\det(A - \lambda \cdot \mathbf{I}_n) = 0.\]
Dette kaldes også for \textbf{karakteristisk polynomium} for $A$ og betegnes $p_{\mathbf{A}}(\lambda)$.

\

Med andre ord, rødderne i følgende polynomie er egenværdierne for $A$.
\[
\det \left(
    \begin{bmatrix}
        a_{11} - \lambda & a_{12} & \cdots & a_{1n} \\
        a_{21} & a_{22} - \lambda & \cdots & a_{2n} \\
        \vdots & \vdots & \ddots & \vdots \\
        a_{n1} & a_{n2} & \cdots & a_{nn} - \lambda 
    \end{bmatrix}\right) = 0
\]

\kilde{Sætning 12.1.1}

\subsubsection{Egenværdi for Lineære afbildinger}
For en lineær afbildning $L: V \to V$ og en basis $\beta$ for $V$, vil en skalar $\lambda \in \mathbb{F}$ være en egenværdi for $L$ hvis og kun hvis
\[
\det(_\beta[L]_{\beta} - \lambda \cdot \mathbf{I}_n) = 0.
\kildetag{Sætning 12.1.2}
\]

\paragraph{Om egenværdier af lineære afbildninger}

\

Den valgte basis er ligegyldig.

\begin{align*}
    p_{_\beta[L]_\beta}(\lambda) &= p_{_\gamma[L]_\gamma}(\lambda) \kildetag{Sætning 12.1.4} \\
    \det(_\beta[L]_\beta) &= \det(_\gamma[L]_\gamma) \kildetag{Korollar 12.1.5}
\end{align*}

Det karakteristiske polynomium for en lineær afbildning $L: V \to V$ betegnes med valgfri basis:
\[
p_L(\lambda) = \det(_\beta[L]_\beta - \lambda \cdot \mathbf{I}_n) \kildetag{Definition 12.1.3}
\]

\subsection{Egenvektorer}
\begin{align*}
    E_\lambda = \{ v \in V | L(v) = \lambda \cdot v \} \kildetag{Sætning 12.2.1} \\
    A_\lambda = \{ v \in V |A \cdot \mathbf{v} = \lambda \cdot \mathbf{v} \} \kildetag{Sætning 12.2.1}
\end{align*}
Disse er underrum i hhv. $V$ og $\mathbb{F}^n$.

\subsection{Egenrum}

Egenrummet for en egenværdi $\lambda$ og matrix $A \in \mathbb{F}^{n \times n}$ betegnes $E_{\lambda}$ og defineres som mængden af alle egenvektorer tilhørende $\lambda$ samt \emph{nulvektoren}. Det kan findes ved at finde kernen af matricen:

\[ 
\begin{bmatrix}
    a_11 - \lambda & a_{12} & \cdots & a_{1n} \\
    a_{21} & a_{22} - \lambda & \cdots & a_{2n} \\
    \vdots & \vdots & \ddots & \vdots \\
    a_{n1} & a_{n2} & \cdots & a_{nn} - \lambda
\end{bmatrix} \kildetag{Lemma 12.2.3}
\]

Se evt. sektion \ref{sec:kerne} om hvordan kernen findes.

\section{Diagonalisering}

\subsection{Kan en matrix diagonaliseres?}
En lineær afbildning $L: V \to V$ eller matrix $A \in \mathbb{F}^{n \times n}$ kan diagonaliseres hvis og kun hvis der findes en basis for $V$ bestående af egenvektorer tilhørende $L$.

Det er det samme som at det karakteristiske polynomium for $L$ kan skrives som
\[
p_L(Z) = (-1)^n \cdot (Z - \lambda_1)^m_1 \cdots (Z - \lambda_r)^{m_r}
\kildetag{Sætning 12.3.3}
\]

Det betyder også at for alle egenværdier $\lambda_i$ gælder
\(
gm(\lambda) = am(\lambda)
\)

Når en afbildnign kan diagonaliseres kan man finde en diagonalmatrix $D$ og en invertibel matrix $Q$ så
\[
L = Q \cdot D \cdot Q^{-1} \kildetag{Lemma 12.4.2}
\]


\subsubsection{Algebraisk multiplicitet}
$am(\lambda)$ er antallet af gange egenværdi $\lambda$ optræder i roden af $p_A(\lambda)$.

F.eks. For $p_A(\lambda) = (\lambda - 2)^3 (\lambda + 1)^2$, er $\am(2) = 3$ og $\am(-1) = 2$. 

\subsubsection{Geometrisk multiplicitet}
\[gm(\lambda) = \dim(E_\lambda)\]

Hvis 
\(
E_{-3} = \spanop \left\{
s \cdot 
\begin{bmatrix}
1 \\ 0 \\ 0
\end{bmatrix} + t \cdot 
\begin{bmatrix}
    0 \\ 1 \\ 0
\end{bmatrix}
\right\}
\)
er $\gm(-3) = 2$.

\subsubsection{Generel sammenhæng}
\[
1 \le \gm(\lambda) \le \am(\lambda) \le \dim(V)
\kildetag{Sætning 12.2.4}
\]

\subsection{Find Diagonalmatrix}
Hvis $A \in \mathbb{F}^{n \times n}$ kan diagonaliseres, kan diagonalmatricen findes ved at indsætte egenværdierne for $A$ på diagonalen i en $n \times n$ matrix $D$.
\[D = 
\begin{bmatrix}
    \lambda_1 & 0 & \cdots & 0 \\
    0 & \lambda_2 & \cdots & 0 \\
    \vdots & \vdots & \ddots & \vdots \\
    0 & 0 & \cdots & \lambda_n
\end{bmatrix}\]

\bemærkBig{Hvis en egenværdi har algebraisk multiplicitet $n$, skal den optræde $n$ gange efter hinanden på diagonalen i $D$.}

\subsection{Find invertibel matrix}
Diagonalmatrix $D$ kan skrives som $D = Q^{-1} \cdot A \cdot Q$, hvor $Q$ er en invertibel matrix bestående af egenvektorerne tilhørende egenværdierne i $D$ som søjler.
\[Q =
\begin{bmatrix}
    | & | & & | \\
    \mathbf{v_1} & \mathbf{v_2} & \cdots & \mathbf{v_n} \\
    | & | & & |
\end{bmatrix}\]

\bemærkBig{Rækkefølgen af egenvektorerne i $Q$ skal svare til rækkefølgen af egenværdierne i $D$.}

\subsection{Koordinater med hensyn til egenvektor, når kun egenværdi er kendt}
Hvis der gives to vektorer $\mathbf{v_1}, \mathbf{v_2} \in V$ og en lineær afbildning $f: V \to V$, og det angives at $\mathbf{v_1}$ er en egenvektor tilhørende egenværdien $\lambda_1$ og $\mathbf{v_2}$ er en egenvektor tilhørende egenværdien $\lambda_2$, kan koordinaterne for \(f(2 \cdot v_1 - 4 \cdot v_2)\) med hensyn til basis \(\beta = \{v_1, v_2\}\) findes som følger:
\[
    (2 \cdot \lambda_1, \quad -4 \cdot \lambda_2)
\]

\subsection{Analyser egenværdier fra polynomium}

For et givent polynomuim \(\lambda^{10} \cdot (\lambda-5)^6 \cdot (\lambda+8)^{10} = 0\), og med forudsætningen at $\lambda_1 < \lambda_2 < \lambda_3$ kan man finde egenværdier samt algebraisk multiplicitet.

Her er Ses at de tre egenværdier, (værdierne det gør polynomiet lig 0) er: $0, 5, -8$, derfor er:
\begin{itemize}
    \item $\lambda_1 = -8$ med $\am(\lambda_1) = 10$
    \item $\lambda_2 = 0$ med $\am(\lambda_2) = 10$
    \item $\lambda_3 = 5$ med $\am(\lambda_3) = 6$
\end{itemize}

\subsection{Find Diagonalmatrix fra egenvektor matrix}

Gives en matrix $A$ \( \begin{bmatrix}
    a_{11} & a_{12} \\
    a_{21} & a_{22}
\end{bmatrix}
\) og en invertibel matrix $Q$ \( \begin{bmatrix}
    v_{11} & v_{12} \\
    v_{21} & v_{22}
\end{bmatrix}
\) bestående af egenvektorerne til $A$ som søjler, kan diagonalmatrix $D$ fines ved at gange hver af søjlerne i $Q$ med $A$, og se hvilket multiplum af søjlen man får.


\subsection{Andet}

For en $2 \times 2$ matrix med kendt determinant og en eigenværdi er den anden egenværdi givet ved
\[\lambda_2 = \frac{\det(A)}{\lambda_1}.\]

\

Med $x$ forskellige egenværdier er der mindst $x$ lineært uafhængige egenvektorer.

\

Egneværdierne for $A^{-1}$ er givet ved $\frac{1}{\lambda_i}$ hvor $\lambda_i$ er egenværdierne for $A$.

\

2 lineært uafhængige egenvektorrum har kun én fælles vektor, nemlig nulvektoren.